\chapter{实验一：熟悉 Oberon-0 语言定义}

\section*{实验目的}

本实验（实验一）旨在深入理解 Oberon-0 语言的定义，为后续词法分析、语法分析和语法制导翻译实验奠定基础。具体目标包括：

\begin{itemize}
    \item 阅读 Oberon-0 语言的 BNF 定义和语义描述，掌握该语言的完整语法规则。
    \item 编写一个词法、语法和语义均正确的 Oberon-0 源程序，覆盖所有语法构造。
    \item 基于变异测试思想，针对每一种可能的词法、语法和语义错误构造变异程序。
    \item 讨论 Oberon-0 语言的特点（保留字与关键字的区别、表达式语法的独特性）。
    \item 分析 Oberon-0 文法的二义性，理解其与其他高级语言在语法设计上的差异。
\end{itemize}

\section*{实验目录与项目组织}

实验一的项目文件均存放于 \texttt{ex1/} 子目录下，各文件说明如下：

\begin{itemize}
    \item \texttt{readme.txt}：自述文件，包含姓名、学号、完成日期及补充说明。
    \item \texttt{testcases/BankSystem.obr}：词法、语法、语义均正确的 Oberon-0 源程序，覆盖了语言的全部语法构造。
    \item \texttt{testcases/BankSystem.001--022}：22 个变异程序（14 个直接构造 + 8 个 derived 派生），分别包含一种词法、语法或语义错误，首行注释标明期望异常类型和变异来源。文件扩展名以数字序号区分。
\end{itemize}

\section{正确的 Oberon-0 源程序}

根据 Oberon-0 语言的 BNF 定义，编写了一个模拟银行账户管理系统的 Oberon-0 源程序 \texttt{BankSystem.obr}，覆盖了 Oberon-0 语言的所有主要语法构造：

\begin{itemize}
    \item \textbf{模块声明}：MODULE BankSystem ... END BankSystem.
    \item \textbf{常量声明}：CONST MaxAccounts = 50; MinBalance = 100;
    \item \textbf{类型声明}：TYPE AccountRec = RECORD ... END; TYPE AccountArray = ARRAY 50 OF AccountRec;
    \item \textbf{变量声明}：VAR accounts, count, i, choice, amount, accId: INTEGER; VAR running, found: BOOLEAN;
    \item \textbf{过程声明}（5 个）：InitAccounts、FindAccount（含 VAR 参数）、Deposit、Withdraw、CreateAccount
    \item \textbf{IF 语句}：含 IF-THEN-ELSE、IF-THEN-ELSIF-THEN-ELSE-END 等完整形式
    \item \textbf{WHILE 循环语句}
    \item \textbf{赋值语句}：含数组下标选择（\texttt{accounts[k].id := 0}）和记录域选择（\texttt{accounts[k].balance := 0}）
    \item \textbf{过程调用}：含实际参数传递
    \item \textbf{预定义过程调用}：Read、Write、WriteLn
    \item \textbf{布尔变量与逻辑运算}：\texttt{~found}、\texttt{(j < count) \& ~found}
\end{itemize}

程序具有实际意义：实现银行账户的创建、存款、取款等基本操作，展示了 Oberon-0 作为教学语言的核心抽象能力。

\section{变异程序}

\begin{table}[htbp]
    \centering
    \begin{tabular}{|c|c|c|c|}
    \hline
    \textbf{变异程序} & \textbf{类别} & \textbf{错误描述} & \textbf{检测结果} \\
    \hline
    BankSystem.001 & 词法 & 非法符号 \$ & IllegalSymbolException $\checkmark$ \\
    BankSystem.002 & 词法 & 标识符超 24 字符 & IllegalIdentifierLengthException $\checkmark$ \\
    BankSystem.003 & 词法 & 非法八进制 08 & IllegalOctalException $\checkmark$ \\
    BankSystem.004 & 词法 & 25id 无空白分隔 & IllegalIntegerException $\checkmark$ \\
    BankSystem.005 & 词法 & 注释未闭合 & MismatchedCommentException $\checkmark$ \\
    BankSystem.014 & 词法 & 整数超12位 & IllegalIntegerRangeException $\checkmark$ \\
    BankSystem.006 & 语法 & 缺左括号 & MissingLeftParenthesisException $\checkmark$ \\
    BankSystem.007 & 语法 & 缺右括号 & MissingRightParenthesisException $\checkmark$ \\
    BankSystem.008 & 语法 & 两表达式间缺运算符 & MissingOperatorException $\checkmark$ \\
    BankSystem.009 & 语法 & 二元运算缺右操作数 & MissingOperandException $\checkmark$ \\
    BankSystem.010 & 语法 & 缺 END & SyntacticException $\checkmark$ \\
    BankSystem.011 & 语义 & INTEGER \& BOOLEAN & TypeMismatchedException $\checkmark$ \\
    BankSystem.012 & 语义 & 参数数目不匹配 & ParameterMismatchedException $\checkmark$ \\
    BankSystem.013 & 语义 & 模块名不匹配 & SemanticException $\checkmark$ \\
    BankSystem.015 & 词法 & 非法符号 @ & IllegalSymbolException $\checkmark$ \\
    BankSystem.022 & 词法 & 非法八进制 09 & IllegalOctalException $\checkmark$ \\
    BankSystem.021 & 词法 & 注释未闭合 & MismatchedCommentException $\checkmark$ \\
    BankSystem.017 & 语法 & 缺左括号 & MissingLeftParenthesisException $\checkmark$ \\
    BankSystem.016 & 语法 & 缺右括号 & MissingRightParenthesisException $\checkmark$ \\
    BankSystem.020 & 语法 & 缺 END & SyntacticException $\checkmark$ \\
    BankSystem.018 & 语义 & \& 用于 INTEGER & TypeMismatchedException $\checkmark$ \\
    BankSystem.019 & 语义 & 参数数目不匹配 & ParameterMismatchedException $\checkmark$ \\
    \hline
    \end{tabular}
    \caption{变异程序及其错误检测结果}
    \label{tab:mutants}
\end{table}

基于变异测试思想，编写了 22 个变异程序（其中 8 个为从 BankSystem.obr 直接最小改动派生的 derived 系列），覆盖 Oberon-0 语言全部 6 种词法错误、4 种语法错误和 2 种语义错误。每个变异程序首行以注释标注期望的错误类型和变异来源。

\section{Oberon-0 语言特点讨论}

\subsection{保留字与关键字的区别}

Oberon-0 语言区别了保留字（Reserved Word）与关键字（Keyword）：

\begin{itemize}
    \item \textbf{保留字}（如 MODULE、END、PROCEDURE、IF、THEN、WHILE、DO、DIV、MOD、OR 等）是语法结构的骨架部分，\textbf{不允许}被程序员用作标识符。在词法分析阶段即被识别为独立 Token 类型。
    \item \textbf{关键字}（如 INTEGER、BOOLEAN、Read、Write、WriteLn、TRUE、FALSE 等）具有预定义语义但\textbf{允许}在嵌套作用域中被重新声明而改变含义。在词法分析阶段被识别为标识符（TK\_IDENT），其特殊语义在语法/语义分析阶段处理。本实验的扫描器将 Read、Write、WriteLn 识别为标识符 Token，由语法分析器根据上下文确定其语义角色。
\end{itemize}

\subsection{表达式语法与 Java/C++ 的差异}

Oberon-0 表达式语法与 Java/C++ 的主要差异：

\begin{enumerate}
    \item \textbf{层次化编码优先级}：Oberon-0 通过 expression $\to$ simple\_expression $\to$ term $\to$ factor 四层产生式自然编码了运算符优先级，无需额外的优先级表。而 C/Java 有 15+ 级优先级，完全依赖优先级约定而非文法结构。
    \item \textbf{逻辑运算符符号不同}：Oberon-0 使用 \&（与）、OR（或）、\~{}（非），不同于 C/Java 的 \&\&、||、!。
    \item \textbf{无短路求值语义}：Oberon-0 未明确规定短路求值，而 C/Java 的 \&\& 和 || 具有短路语义。
    \item \textbf{运算符集合少}：不支持位运算、自增自减、三目运算符、复合赋值。
\end{enumerate}

\section{Oberon-0 文法二义性讨论}

经分析，Oberon-0 的 BNF 定义\textbf{不存在二义性}：

\begin{enumerate}
    \item \textbf{悬挂 else 问题已被消除}：Oberon-0 的 IF 语句使用 IF-THEN-ELSIF-ELSE-\textbf{END} 闭合结构，每个 IF 必须有显式的 END 结束标记。这从根本上避免了 Pascal/C 等语言中 else 与哪个 if 匹配的经典二义性。Wirth 在其语言设计中一贯坚持这一理念。
    \item \textbf{运算符优先级内置于文法层次}：四层表达式文法结构（expression $\to$ simple\_expression $\to$ term $\to$ factor）已明确编码了关系运算、加减运算、乘除运算的优先级关系，无需额外的优先级和结合性规则。
    \item \textbf{运算符重载通过位置消歧}：一元 \texttt{-} 出现在 simple\_expression 起始的 \texttt{["+" | "-"]} 中，二元 \texttt{-} 出现在 \texttt{("+" | "-" | "OR")} 中，语法位置自然区分。
\end{enumerate}

\section{自我评分参照}

参照实验一评价标准，各项完成情况如下：

\begin{table}[htbp]
    \centering
    \begin{tabular}{|p{4.5cm}|c|c|p{4.5cm}|}
    \hline
    \textbf{项目} & \textbf{权重(\%)} & \textbf{对应内容} \\
    \hline
    1、正确的 Oberon-0 源程序 & 20 & 3.1节 BankSystem.obr \\
    2、变异程序 & 40 & 3.2节 22个变异程序（含8个derived） \\
    3、语言特点讨论 & 10 & 3.3节 保留字/关键字、表达式差异 \\
    4、文法二义性讨论 & 20 & 3.4节 无二义性证明 \\
    5、文档组织及 readme & 10 & ex1/readme.txt \\
    \hline
    \end{tabular}
    \caption{实验一自我评分参照}
    \label{tab:self-eval-ex1}
\end{table}

\section{实验心得体会}

通过实验一，我对 Oberon-0 语言有了系统而深入的理解。编写覆盖全部语法构造的正确源程序 BankSystem.obr 的过程，让我逐一梳理了 Oberon-0 的模块结构、声明体系、控制流构造和表达式语法。在设计变异程序时，我深刻体会到变异测试"最小改动、精确触发"的思想——每个变异仅改变程序的一处细节，便能准确揭露特定类别的错误检测需求。

Oberon-0 语法设计的简洁与严谨给我留下了深刻印象。Wirth 通过 IF-END 显式闭合结构从根本上消除了悬挂 else 问题，通过四层表达式文法将运算符优先级自然编码在语法中——这些设计使得 Oberon-0 成为一门分析友好、教学价值极高的语言。与 Java/C++ 等工业语言的对比也让我更清楚地看到，语言的"简洁性"与"表达能力"之间始终存在权衡。

